\section{Cuadrado mágico cuántico no semiclásico en una esquina APP}
\label{sec:cpe-qms-no-semiclasico}

En \(\mathbb C^4\), con base \(e_1,e_2,e_3,e_4\), considérese el arreglo de
vectores
\[
\begin{array}{cccc}
e_1&e_2&e_3&e_4\\
e_2&e_1&e_4&e_3\\
(e_3+e_4)/\sqrt2&(e_3-e_4)/\sqrt2&(e_1+e_2)/\sqrt2&(e_1-e_2)/\sqrt2\\
(e_3-e_4)/\sqrt2&(e_3+e_4)/\sqrt2&(e_1-e_2)/\sqrt2&(e_1+e_2)/\sqrt2.
\end{array}
\]
Sea \(P_{ij}\) el proyector de rango uno sobre el vector situado en la
celda \((i,j)\). Cada fila y cada columna es una base ortonormal, luego
\begin{equation}
 \sum_jP_{ij}=I_4,
 \qquad \sum_iP_{ij}=I_4.
 \label{eq:cpe-qms4}
\end{equation}
Para realizarlo en la fibra APP de dimensión \(9\), se fija
\begin{equation}
 A_{ij}:=P_{ij}\oplus\frac14I_5\in M_9(\mathbb C).
 \label{eq:cpe-qms9}
\end{equation}

\begin{theorem}[No semiclasicidad certificada]
\label{thm:cpe-qms-no-semiclasico}
La matriz \(A=(A_{ij})_{i,j=1}^4\) es un cuadrado mágico cuántico y no admite
una descomposición semiclásica
\[
 A_{ij}=\sum_{\pi:\,\pi(i)=j}Q_\pi,
 \qquad Q_\pi\succeq0,
 \qquad\sum_\pi Q_\pi=I_9.
\]
\end{theorem}

\begin{proof}
Las identidades de fila y columna se siguen de
\eqref{eq:cpe-qms4}; cuatro copias de \(I_5/4\) suman \(I_5\). Si existiera
una descomposición semiclásica, su compresión a la esquina \(\mathbb C^4\)
sería una medición conjunta de las PVM definidas por las filas
\((P_{ij})_j\). Las mediciones proyectivas conjuntamente medibles conmutan.
Sin embargo,
\[
 \left[
 |e_1\rangle\langle e_1|,
 \left|\frac{e_1+e_2}{\sqrt2}\right\rangle
 \left\langle\frac{e_1+e_2}{\sqrt2}\right|
 \right]
 =\frac12
 \begin{pmatrix}0&1\\-1&0\end{pmatrix}
 \oplus0_2\ne0.
\]
La contradicción prueba la no semiclasicidad. La compresión proporciona, a
la vez, un testigo analítico de inviabilidad del problema semidefinido.
\end{proof}

Cada fila define un canal UCP
\[
 \Phi_i:\ell^\infty_4\to M_9,
 \qquad \Phi_i(\delta_j)=A_{ij},
\]
y cada columna define otro. Así, las normalizaciones mágicas se traducen a
mapas completamente positivos con dominio y codominio explícitos; no se usa
un símbolo \(J(\operatorname{id})\) sin haber fijado esos tipos.

\section{Sistema cónico de rutas y universalidad genealógica interna}
\label{sec:cpe-s-system-universalidad}

Para cada perfil finito \(c\), sea \(G_c\) el grupoide finito de rutas TPK
con ese perfil y sea \(A(c)=C^*(G_c)\). En nivel matricial \(s\), se toma el
cono
\[
 \mathcal C_s(c)=M_s(A(c))_+.
\]
Una ruta \(r:c\to d\), representada por la isometría parcial \(v_r\),
induce el mapa CP
\[
 T_r(a)=v_r^*av_r.
\]

\begin{theorem}[Sistema cónico functorial de las rutas TPK]
\label{thm:cpe-s-system-tpk}
Las asignaciones \(c\mapsto(\mathcal C_s(c))_{s\ge1}\) y
\(r\mapsto T_r\) respetan identidad, composición, amplificación matricial y
truncamiento genealógico. La categoría CPM generada por las fibras finitas
es monoidal compacta: el dual de una fibra es su espacio conjugado y las
coevaluaciones y evaluaciones son las copas y tapas canónicas. El límite
proyectivo de estos sistemas conserva todos los mapas CP compatibles.
\end{theorem}

\begin{proof}
La positividad completa de \(a\mapsto v_r^*av_r\) es inmediata de su forma
de Kraus. Para rutas composables,
\(v_{r_2r_1}=v_{r_2}v_{r_1}\), luego
\(T_{r_2r_1}=T_{r_1}T_{r_2}\) con la varianza declarada. Las inclusiones y
esperanzas condicionales de los niveles finitos son UCP y forman cuadrados
conmutativos. La compactación de CPM en dimensión finita es la construcción
estándar por espacios conjugados; la compatibilidad de truncación permite
pasar al límite coordenada a coordenada.
\end{proof}

Sea \(\mathsf P_{\rm HMT}\) el conjunto numerable de códigos de rutas
finitas y sea \(\mathsf C_{\rm HMT}\) el espacio de estados con su contexto
completo. La evaluación
\[
 \operatorname{Ev}(\ulcorner r\urcorner,x)=T_r(x)
\]
conserva origen, destino, hoja, orientación, acarreo, frontera y palabra de
ruta.

\begin{corollary}[Universalidad genealógicamente fiel en el dominio HMT]
\label{cor:cpe-universalidad-genealogica}
El evaluador \(\operatorname{Ev}\) es universal para la clase declarada de
transportes TPK finitos y sus límites compatibles: todo transporte de esa
clase posee un código, la composición se realiza concatenando códigos y la
reducción conserva la genealogía completa. Esta universalidad interna no se
extiende por decreto a una clase externa de objetivos no tipada.
\end{corollary}

\section{Identificación de la memoria de Stinespring con el registro local TPK}
\label{sec:cpe-stinespring-tpk}

En la esperanza nonádica, la dilatación de Stinespring conserva en una base
\(\{e_a\}_{a=0}^8\) el dígito descartado. Para un estado TPK \(x\), sea
\(\lambda(x,a)\) la etiqueta completa de la continuación por el dígito
\(a\). Se define la isometría
\[
 W_x:e_a\longmapsto e_{\lambda(x,a)}
\]
desde la memoria mínima de Stinespring hacia el registro local TPK\@. Sea
\(F_xe_{\lambda(x,a)}=e_a\) el funtor de olvido restringido a su imagen.

\Needspace{12\baselineskip}
\begin{theorem}[La memoria TPK refina la memoria mínima de Stinespring]
\label{thm:cpe-stinespring-tpk}
Para todo \(x\),
\[
 W_x^*W_x=I_9,
 \qquad F_xW_x=I_9.
\]
La familia \((W_x)_x\) entrelaza la inclusión de un nuevo dígito y los
truncamientos. En profundidad \(m\), el tensor de estas isometrías identifica
la ancilla \((\mathbb C^9)^{\otimes m}\) con el subregistro TPK que conserva
la sucesión de extensiones; las coordenadas adicionales de hoja, acarreo,
frontera y ruta explican por qué la memoria TPK es un refinamiento y no una
segunda copia de la ancilla.
\end{theorem}

\begin{proof}
Continuaciones con dígitos distintos poseen etiquetas ortogonales, por lo
que \(W_x\) es isométrica. La definición de \(F_x\) prueba la segunda
identidad. La composición de extensiones concatena las etiquetas en el mismo
orden que el producto tensorial de las ancillas; de ahí la naturalidad bajo
truncamiento y la afirmación a profundidad arbitraria.
\end{proof}
